\section{Introduction}
\label{a:sec:intro}

For the binary Galton--Watson process with division probability $p\in[0,\tfrac12)$
the survival probability decays geometrically, $S_n\sim A(p)(2p)^n$, and the mean
population conditional on non-extinction converges to $1/A(p)$. That much was
established in \ChM. Everything the limiting conditional law does --- its mean, its
variance, its dependence on how close the process sits to criticality --- passes
through the single constant
\begin{equation}
\label{a:eq:Apdef}
A(p) = \lim_{n\to\infty}\frac{S_n}{(2p)^n},
\qquad p\in[0,\tfrac12) .
\end{equation}
This chapter is about that constant: what it is, how to compute it, why no
elementary formula for it has been found despite a determined search, and exactly
what can and cannot be proved about the failure of that search. The results are
original to this thesis except where attributed.

Approximate values are easy enough to obtain. Iterate
$S_{n+1}=2pS_n-pS_n^2$ until $S_n$ is very nearly stationary, then form
$S_n/(2p)^n$. Repeating over a range of $p$ traces the curve of
\cref{a:fig:dtct}, and this works well over most of the interval. It fails as
$p\to\tfrac12^-$, because the number of generations needed for $S_n$ to reach
near-stationarity grows without bound --- from around $10^2$ at $p=0.2$ to nearly
$5\times10^5$ at $p=0.4999$.

After a short recapitulation of the setup inherited
from \ChM{} (\cref{a:sec:setup}), two exact representations come first. Rewriting
the survival recursion as a telescoping product (\cref{a:sec:product}) proves that
the limit \cref{a:eq:Apdef} exists and is strictly positive; converting that
product into a series (\cref{a:sec:series}) isolates the geometric decay in a
factor carrying no dynamics at all, so that two-sided bounds can be read off by
inspection. Those bounds trap $A(p)$ between half the continuous-time constant
$\Ac(p)$ and a companion upper bound, and a parity argument caps it at $\tfrac12$.
The same series then yields the near-critical behaviour
$A(p) = 2\varepsilon[1+O(\varepsilon\log(1/\varepsilon))]$, with
$\varepsilon = 1-2p$ (\cref{a:sec:nearcritical}), locating the logarithm that makes
convergence to the leading term slow, and with it the exact relationship between
$A$ and $\Ac$ (\cref{a:sec:compare}).

What none of that yields is a formula. \Cref{a:sec:GA} records the search for one
and the manner of its failure. The explanation is dynamical: under the
substitution $S_n = 2w_n$ the survival
recursion becomes the logistic map, and $A(p)$ is twice a value of that map's
Koenigs linearising function (\cref{a:sec:koenigs}); for every multiplier in the
subcritical window that function is hypertranscendental (\cref{a:sec:ht}), so no
elementary expression in its own variable can be evaluated to produce $A(p)$.
\Cref{a:sec:scope} delimits what this does and does not settle and reports the
inverse-symbolic evidence bearing on what it leaves open, and
\cref{a:sec:practical,a:sec:conclusion} draw the practical and the residual
conclusions.
